\section{Continuity}

Everything else in the bridge depends on this rung. Until it runs, the claim that the
coarse-grained face of the base is a conserved-current field theory is a
well-posed conjecture. After it runs, the central identification is a measured
fact on the lattice. Coarse-grain the conserved
tone-current over growing blocks and its discrete divergence is exactly zero at
every block scale. The charge change inside each block equals minus the net flux
across its boundary, with no residual. That is the discrete statement of
$\nabla\!\cdot J = 0$, the one law a conserved-current program posits, and from
which it forces everything else \cite{herbert2025chronoflux}.

\begin{finding}
On the flat torus (\texttt{E-GRV-0039}) the per-block residual is an integer-exact zero
at block scales $1$, $2$, and $4$, with no tolerance allowed. The conserving
rule is divergence-free at every coarse-graining. The lossy control injects a
per-cell sink and its residual is $2752$, $704$, and $192$ at the same scales,
nonzero everywhere and never shrinking. On the genuinely curved $\{3,4,3,4\}$
mesh (\texttt{E-GRV-0042}) the per-region residual is again integer-exact zero at every
ball size and every beat over thousands of curved cells, with global
conservation and exact reversibility, while the lossy rule breaks it by exactly
the charge it destroys inside each ball. Tier F, \textbf{Derived}.
\end{finding}

The argument is worth walking through, because the identity is so clean it can
look like an accident. Every dock carries one tone per direction. From these,
build two plain things. The density is the signed sum of the tones, one number
saying how much charge sits at the dock. The current is each tone weighted by
its direction and summed, one arrow saying which way the charge flows.
That pair is the tone-current, built from what is there. Pick
any patch of docks and call it a block. Its boundary is the set of direction
links that cross from inside to outside, and each crossing link carries one unit
of charge stepping over the edge in one beat.

Now ask why the charge inside changes. The rule does two things and neither
makes nor destroys charge. Collide rotates tone among a dock's head-on
direction pairs and leaves the dock's total untouched. Stream slides each tone to
the neighbor it points at, so a charge that leaves one dock arrives whole at the
next. Nothing evaporates and nothing appears. The only way the charge inside a
block can change over a beat is for some charge to physically walk across the
boundary. The change inside equals the net amount that crossed the edge exactly,
by the way the rule is built. That sentence is the discrete
divergence theorem.

Block-averaging carries the per-block balance up to the field equation
untouched. Average over blocks of growing size and the per-dock jitter blurs
while the density and current settle into smooth fields. The same sum rule,
holding at every block and every scale, is exactly the statement that the rate
of change of density plus the divergence of the current is zero. That is
$\nabla\!\cdot J = 0$, the same sentence read at a coarser grain. The lossy
control is the load-bearing check. Its sink lets charge vanish from a block
without crossing the boundary, so the books do not balance and the residual stays
nonzero at every scale. Closure comes from conservation. It is not a byproduct
of the smoothing that any coarse-graining produces.

So there are not two laws. A continuum program posits continuity at the field
scale and forces the rest from it. The base never posits it. The base conserves
charge dock by dock, and block-averaging cannot break a balance the rule keeps
exactly at every step. Charge conservation at the cell scale is continuity at
the field scale, one law seen at two grains. The remaining piece is the
explicit map from the lattice sum rule to a single named continuum field with
its flow direction and its stress. The base is literally a lattice gas, the
canonical engine that produces emergent hydrodynamics from one conserved local
rule \cite{frisch1986, dhumieres1986}, and the curved run so far streams along
the Coxeter generators rather than the full 24-direction coin, so carrying the
complete direction set onto the curved mesh is the next step. The conservation
is exact because the rule is reversible \cite{kari2005, fredkin1982}, and the
recovery of a smooth four-dimensional continuum from discrete building blocks is
the same destination the causal-triangulation program reaches by a different
road \cite{ambjorn2004, cluesurf2026repo}.
